\documentclass[10pt,a4paper]{amsart}
\usepackage[margin=19mm]{geometry}
\usepackage{amsmath,amssymb,mathtools}
\usepackage[scr=rsfs,cal=euler]{mathalfa}
\usepackage{xfrac}
\usepackage[unicode,hidelinks,bookmarks=false]{hyperref}
\newcommand{\ie}{i.e.\ }
\newcommand{\cf}{cf.\ }
\newcommand{\resp}{respectively\ }
\newcommand{\Subsection}{\subsection}
\providecommand{\nobreakdash}{\nobreak\hbox{-}}
\setlength{\emergencystretch}{2em}
\multlinegap0pt
\usepackage{mathtools}
\usepackage{amssymb}
\usepackage{float}
\usepackage{xcolor}
\newtheorem{lemma}{Lemma}
\newcommand{\M}{\mathcal{M}}
\newcommand{\T}{\mathcal{T}}
\title{On the maximum number of vectors in $\{0,\pm1\}^n$ with forbidden inner products}
\author{Ilya Lobatskii, Yakov Shubin}
\date{Source: arXiv:2606.12178v1 (CC BY 4.0)}
\begin{document}
\maketitle

\noindent\emph{Source and licence.} On the maximum number of vectors in $\{0,\pm1\}^n$ with forbidden inner products. Version arXiv:2606.12178v1 (CC BY 4.0), distributed by arXiv under the Creative Commons Attribution 4.0 International licence, http://creativecommons.org/licenses/by/4.0/, as recorded in the arXiv licence metadata for this exact version and shown on the abstract page https://arxiv.org/abs/2606.12178v1. Full licence notice: \texttt{licenses/arxiv-2606.12178-CC-BY-4.0.txt}.\par\smallskip

\noindent\textit{Editorial scope.} Counted unit: the article's lemma thm3 (source0.tex lines 338--340). The article's complete original proof of it is delivered with it (source0.tex lines 342--374, one \texttt{proof} environment of 33 lines). No mathematics is written, repaired or re-derived: the statement and the proof are the article's own lines, in the article's own notation, and no retained line is substituted or re-flowed.  Retained uncounted as the source-local context the proof needs: lines 306--308.  Nothing was omitted from any retained span, so the package carries no omission record.  No retained line contains a citation, so the wrapper carries no bibliography and no bibliography entry.  No retained line contains a figure environment, an image-inclusion command or drawing code, so no image is delivered and the package declares no asset dependency.  Editorial formatting only, in addition: an AMS \texttt{amsart} preamble that restates the article's own declarations and macros used by the retained lines, namely the environments \texttt{lemma} on separate counters, together with the standard packages those lines need.  The retained environments print in this document as Lemma 1 in order of appearance, whereas the article prints the same items with the numbering of its own full document.  The complete native source of the article as distributed in the arXiv e-print for this exact version is bundled unchanged as \texttt{sources/202604/source0.tex}.
\begin{lemma}\label{mainlem}
    If $|\M| \geqslant 8n-48$, then $k+l+m<30$.
\end{lemma}

\begin{lemma}\label{thm3}
    There exists $n_0$ such that, if $n>n_0$ and $|\M| \geqslant 8n-48$, then there are two disjoint serpentaria.
\end{lemma}

\begin{proof}
    For a vector $v$ we write $p_{pm}(v)=p(v)\cup\{-a:a\in p(v)\}$. Define $p_{pm}(F)$ analogously for a set of vectors $F$.
    Define the sets $H_1,H_2,H_3,H_4 \subset P$ as follows:
    \[
    H_1 \coloneqq p_{pm}(\{h \bigm| \exists i \in [k], h   \text{ is the head of } S_i\}),
    \]
    \[
    H_2 \coloneqq \bigcup_{i\in [l]} p_{pm}(C_i),
    \]
    \[
    H_3 \coloneqq \bigcup_{v \in \T} p_{pm}(v),
    \]
    \[
    H_4 \coloneqq \bigcup_{i \in [k], |S_i| \leqslant 3} p_{pm}(S_i),
    \]
    \[
    H \coloneqq P \setminus (H_1 \cup H_2 \cup H_3 \cup H_4).
    \]
    From the definitions it follows that $|H_1| \leqslant 6k$, $|H_2| \leqslant 10l$, $|H_3| \leqslant 8m$, and $|H_4| \leqslant 12k$. By Lemma~\ref{mainlem}, we have $k+l+m<30$. Therefore
    $$|H_1|+|H_2|+|H_3|+|H_4| \leqslant 18(k+l+m) \leqslant 18\cdot29=522.$$
    Hence
    $$|H| \geqslant |P|-|H_1|-|H_2|-|H_3|-|H_4| \geqslant 2n-522.$$

    We prove that for sufficiently large $n$ there exists a position $a \in H$ such that $f(a)=f(-a)=4$. We argue by contradiction. Note that if $a \in H$, then $-a \in H$. Let $n>309$. Then $|H|>2\cdot309-522=96$. We have
    \[
    \sum_{a\in P} f(a) = \sum_{a\in H} f(a) + \sum_{a\in P \setminus H} f(a)
    \leqslant
    4|P| - \frac{1}{2}|H| < 8n-48 \implies |\M|<8n-48.
    \]
    This is a contradiction. Now fix $a\in H$ for which $f(a)=f(-a)=4$. Since $a\notin H_2\cup H_3\cup H_4$, every class that uses $a$ is a snake of cardinality at least 4. Since $a\notin H_1$, the position $a$ does not lie in the head of such a snake. We show that the snakes containing the position $a$ are pairwise compatible. Consider any two distinct such snakes $S_1,S_2$ with heads $h_1,h_2$, respectively. Then
    $$h_1+a\in S_1, \quad h_2+a\in S_2 \implies (h_1+a,h_2+a)\leqslant 0 \implies (h_1,h_2)\leqslant -1,$$
    as required. Thus, the snakes containing the position $a$ form a serpentarium. The same argument applies to $-a$. It remains only to observe that these serpentaria cannot have common snakes, since a snake cannot use both positions $a$ and $-a$ at the same time. Hence we have found two disjoint serpentaria, and the lemma is proved.
\end{proof}

\end{document}
